% titel: Blatt 6, A3: Polarkoordinaten
% Teil II: der durchgerechnete Weg.

\subsection*{Gegeben und gesucht}
Polarkoordinaten mehrerer Zahlen (a), Potenz $(1-i\sqrt3)^{10}$ (b),
Lösungen von $z^4=-1$ (c), Drehung um $\pi/3$ von $z=\sqrt3+i$ (d).

\subsection*{a) Polarkoordinaten}
Die Zahlen sind reell an den Achsen bzw. auf einfachen Linien:
\begin{itemize}
  \item $-5$: $r=5$, $\varphi=\pi$;
  \item $-3i+3=3-3i$: $r=\sqrt{9+9}=3\sqrt2$,
        $\varphi=-\frac\pi4$;
  \item $\frac1i=-i$: $r=1$, $\varphi=-\frac\pi2$;
  \item $1+i\tan\alpha$: $r=\sqrt{1+\tan^2\alpha}=\sec\alpha$,
        $\varphi=\alpha$.
\end{itemize}

\subsection*{b) $(1-i\sqrt3)^{10}$}
Schreibe $1-i\sqrt3=2\mathrm e^{-i\pi/3}$ (Radius $2$, Winkel $-\pi/3$, im
vierten Quadranten). Dann
\begin{equation*}
  (1-i\sqrt3)^{10}=2^{10}\,\mathrm e^{-i\frac{10\pi}{3}}
  =1024\,\mathrm e^{-i\frac{10\pi}{3}}.
\end{equation*}
Den Winkel modulo $2\pi$ reduzieren:
$-10\pi/3\equiv -10\pi/3+2\cdot2\pi=\frac{2\pi}{3}$ (mod $2\pi$):
\begin{equation*}
  1024\,\mathrm e^{i\frac{2\pi}{3}}
  =1024\Bigl(\cos\frac{2\pi}{3}+i\sin\frac{2\pi}{3}\Bigr)
  =1024\Bigl(-\frac12+i\frac{\sqrt3}{2}\Bigr)
  =-512+512\sqrt3\,i.
\end{equation*}

\subsection*{c) Lösungen von $z^4=-1$}
$-1=\mathrm e^{i\pi}$. Mit $r=1$ und der Wurzelformel:
$z_k=e^{i(\pi+2\pi k)/4}$ für $k=0,1,2,3$:
\begin{equation*}
  z_0=e^{i\pi/4},\ \ z_1=e^{i3\pi/4},\ \ z_2=e^{i5\pi/4},\ \ z_3=e^{i7\pi/4},
\end{equation*}
d.h. $z=\pm\frac{1+i}{\sqrt2}$ und $z=\pm\frac{1-i}{\sqrt2}$: die vier
Punkte auf dem Einheitskreis bei $45^\circ,135^\circ,225^\circ,315^\circ$.

\subsection*{d) Drehung von $z=\sqrt3+i$}
$z=\sqrt3+i$ hat $r=\sqrt{3+1}=2$ und $\varphi=\arctan\frac1{\sqrt3}
=\frac\pi6$. Drehung um $\pi/3$ ergibt
\begin{equation*}
  z'=2\,\mathrm e^{i(\frac\pi6+\frac\pi3)}=2\,\mathrm e^{i\frac\pi2}=2i.
\end{equation*}

\begin{ergebnis}
\begin{equation*}
  z'=2i\ \text{(d)};\qquad z^4=-1\ \Leftrightarrow\
  z=\pm\frac{1\pm i}{\sqrt2}.
\end{equation*}
\end{ergebnis}

\begin{kontrolle}
\textbf{Probe b):} $(1-i\sqrt3)^2=-2-2i\sqrt3$; daraus
$(\cdot)^5$ etc.\ liefert dieselbe Zahl, oder direkt: $r^{10}=2^{10}=1024$,
Winkel $10(-\pi/3)=-10\pi/3\equiv2\pi/3$ -- Wert $-512+512\sqrt3 i$
passt mit dem Quadranten ($x<0,y>0$, zweiter Quadrant, Winkel $2\pi/3$).

\textbf{Probe d):} $z'=2i$, Länge bleibt $2$, Winkel $\pi/2$ -- richtig
gedreht um $+\pi/3$.
\end{kontrolle}

\begin{merksatz}
Potenzen und Wurzeln immer in Polarkoordinaten mit der Moivre-Formel rechnen:
$z^n=r^n e^{in\varphi}$, $z^{1/n}=r^{1/n}e^{i(\varphi+2\pi k)/n}$. Winkel
modulo $2\pi$ reduzieren. Drehung = Multiplikation mit $e^{i\theta}$.
\end{merksatz}
